\begin{onehalfspace}

\section{Il Problema della Ricerca Non Strutturata}

Dato un database\footnote{Il termine \textit{database} è qui usato in senso astratto e generale: non si intende necessariamente una base di dati relazionale, ma qualsiasi insieme di $N$ elementi su cui sia definita una funzione oracolo $f$. L'algoritmo di Grover è applicabile a qualsiasi problema in cui la verifica di una soluzione sia computazionalmente semplice, mentre la ricerca non lo è.} di $N$ elementi, la ricerca di un elemento che soddisfi una certa proprietà richiede, nel caso classico peggiore, $O(N)$ interrogazioni: è necessario esaminare gli elementi uno per uno. Se la struttura del database non consente alcuna indicizzazione o ordinamento, questo limite è ottimale per qualsiasi algoritmo deterministico o probabilistico classico.

Formalmente, il problema è definito come segue. Sia $f:\{0,1,\ldots,N-1\} \to \{0,1\}$ una funzione \textit{oracolo} tale che:
\begin{equation}
    f(x) = \begin{cases} 1 & \text{se } x \text{ è la soluzione (o una delle soluzioni)} \\ 0 & \text{altrimenti} \end{cases}
\end{equation}

Si assume che esista esattamente una soluzione $x^*$ con $f(x^*) = 1$ (si generalizzerà al caso di $M$ soluzioni). L'obiettivo è trovare $x^*$ con il minor numero possibile di interrogazioni a $f$.

Nel 1996, Lov Grover\footnote{\textbf{Lov Kumar Grover} (nato nel 1961 in India) è un informatico teorico che lavorava, come Shor, presso i Bell Labs di AT\&T. La coincidenza è notevole: i due algoritmi quantistici più importanti della storia furono sviluppati nello stesso laboratorio a pochi anni di distanza.} dimostrò che un algoritmo quantistico può trovare $x^*$ con $O(\sqrt{N})$ interrogazioni all'oracolo \cite{grover1996}, ottenendo uno \textbf{speedup quadratico} rispetto al caso classico.

\section{L'Oracolo Quantistico}

L'oracolo classico $f$ viene codificato in un operatore quantistico unitario $O_f$. Esistono due formulazioni equivalenti:

\subsection{Oracolo Computazionale}
\begin{equation}
    O_f\ket{x}\ket{b} = \ket{x}\ket{b \oplus f(x)}
\end{equation}
dove $\ket{b}$ è un qubit ancilla. Questa formulazione è analoga all'oracolo classico.

\subsection{Oracolo di Fase (Phase Oracle)}
Inizializzando il qubit ancilla nello stato $\ket{-} = (\ket{0}-\ket{1})/\sqrt{2}$ e sfruttando il \textit{phase kickback}:
\begin{equation}
    O_f\ket{x}\ket{-} = (-1)^{f(x)}\ket{x}\ket{-}
\end{equation}

Il qubit ancilla rimane invariato, ma lo stato del registro principale acquisisce una fase $-1$ se e solo se $x$ è la soluzione:
\begin{equation}
    O_f\ket{x} = (-1)^{f(x)}\ket{x}
\end{equation}

Questa formulazione è quella utilizzata nell'algoritmo di Grover: l'oracolo \textbf{marca} la soluzione invertendone il segno dell'ampiezza.

\section{Il Meccanismo di Amplitude Amplification}

L'algoritmo di Grover si basa su un operatore composto, il \textbf{Grover Iterator} $G$, che viene applicato ripetutamente per amplificare progressivamente l'ampiezza dello stato soluzione.

\subsection{Inizializzazione}

Si parte da $n = \lceil \log_2 N \rceil$ qubit inizializzati in $\ket{0}^{\otimes n}$. Si applica la trasformata di Hadamard per creare la sovrapposizione uniforme:
\begin{equation}
    \ket{s} = H^{\otimes n}\ket{0}^{\otimes n} = \frac{1}{\sqrt{N}}\sum_{x=0}^{N-1}\ket{x}
\end{equation}

Nello stato $\ket{s}$, ogni stato base ha ampiezza $1/\sqrt{N}$, corrispondente a una probabilità di misura $1/N$ per ciascuno.

\subsection{Il Grover Iterator}

Il Grover Iterator è definito come:
\begin{equation}
    G = D \cdot O_f
\end{equation}

dove:
\begin{itemize}
    \item $O_f$ è l'oracolo di fase (inverte il segno dell'ampiezza della soluzione);
    \item $D$ è il \textbf{Diffusion Operator} (o Grover Diffusion Operator):
    \begin{equation}
        D = 2\ket{s}\bra{s} - I = H^{\otimes n}(2\ket{0}\bra{0} - I)H^{\otimes n}
    \end{equation}
\end{itemize}

Il Diffusion Operator realizza una \textbf{riflessione attorno alla media}: ogni ampiezza $\alpha_x$ viene trasformata in $2\bar{\alpha} - \alpha_x$, dove $\bar{\alpha}$ è la media delle ampiezze. Questa operazione amplifica le ampiezze superiori alla media e attenua quelle inferiori.

\subsection{Interpretazione Geometrica}

L'algoritmo di Grover ammette un'elegante interpretazione geometrica. Si definiscono i due stati ortonormali:
\begin{align}
    \ket{w}  &= \ket{x^*} \quad \text{(la soluzione)} \\
    \ket{s'} &= \frac{1}{\sqrt{N-1}}\sum_{x \neq x^*}\ket{x} \quad \text{(tutti gli altri stati)}
\end{align}

Lo stato iniziale $\ket{s}$ si decompone come:
\begin{equation}
    \ket{s} = \sin\theta\ket{w} + \cos\theta\ket{s'}, \qquad \theta = \arcsin\!\left(\frac{1}{\sqrt{N}}\right)
\end{equation}

Il Grover Iterator $G$ equivale a una rotazione di angolo $2\theta$ nel piano bidimensionale generato da $\ket{w}$ e $\ket{s'}$:
\begin{equation}
    G^k\ket{s} = \sin\!\left((2k+1)\theta\right)\ket{w} + \cos\!\left((2k+1)\theta\right)\ket{s'}
\end{equation}

Dopo $k$ iterazioni, la probabilità di misurare la soluzione è:
\begin{equation}
    P(x^*) = \sin^2\!\left((2k+1)\theta\right)
\end{equation}

\section{Numero Ottimale di Iterazioni}

La probabilità di successo è massimizzata quando $(2k+1)\theta \approx \pi/2$, ovvero:
\begin{equation}
    k_{\text{opt}} = \left\lfloor \frac{\pi}{4\theta} \right\rfloor \approx \frac{\pi}{4}\sqrt{N}
\end{equation}

poiché per $N$ grande $\theta \approx 1/\sqrt{N}$. Il numero ottimale di iterazioni è quindi $O(\sqrt{N})$.

\textbf{Attenzione:} Applicare troppe iterazioni è dannoso. Superato il massimo, la probabilità decresce: il sistema ``supera'' la soluzione. In un caso ideale, la probabilità di successo è:
\begin{equation}
    P_{\text{success}} = \sin^2\!\left(\left(2k_{\text{opt}}+1\right)\theta\right) \geq 1 - \frac{1}{N}
\end{equation}

che tende a 1 per $N$ grande.

\subsection{Caso con $M$ Soluzioni}

Se il database contiene $M$ soluzioni, l'angolo iniziale diventa $\theta = \arcsin(\sqrt{M/N})$ e il numero ottimale di iterazioni diventa:
\begin{equation}
    k_{\text{opt}} \approx \frac{\pi}{4}\sqrt{\frac{N}{M}}
\end{equation}

La complessità rimane $O(\sqrt{N/M})$, che converge a $O(1)$ quando $M \sim N$.

\section{Schema dell'Algoritmo}

\begin{algorithm}[H]
\caption{Algoritmo di Grover}
\begin{algorithmic}[1]
\Require Oracolo $O_f$, numero di elementi $N = 2^n$, numero di soluzioni $M$
\Ensure Soluzione $x^*$ con alta probabilità
\State Inizializzare $n$ qubit: $\ket{\psi_0} = \ket{0}^{\otimes n}$
\State Applicare $H^{\otimes n}$: $\ket{s} = \frac{1}{\sqrt{N}}\sum_{x=0}^{N-1}\ket{x}$
\State Calcolare $k_{\text{opt}} = \lfloor\frac{\pi}{4}\sqrt{N/M}\rfloor$
\For{$k = 1$ \textbf{to} $k_{\text{opt}}$}
    \State Applicare $O_f$: \quad $\ket{x} \mapsto (-1)^{f(x)}\ket{x}$ \Comment{Marca la soluzione}
    \State Applicare $D = H^{\otimes n}(2\ket{0}\bra{0}-I)H^{\otimes n}$ \Comment{Riflessione attorno alla media}
\EndFor
\State Misurare in base computazionale
\State \Return il valore misurato
\end{algorithmic}
\end{algorithm}

\section{Ottimalità e Lower Bound}

L'algoritmo di Grover è \textbf{ottimale}: non esiste alcun algoritmo quantistico che risolva il problema della ricerca non strutturata in meno di $\Omega(\sqrt{N})$ interrogazioni all'oracolo. Questo risultato è stato dimostrato da Bennett et al. (1997) \cite{bennett1997} utilizzando il metodo dei polinomi e, in forma più generale, da Zalka (1999) \cite{zalka1999}.

La dimostrazione si basa sul fatto che le ampiezze delle soluzioni, dopo $k$ interrogazioni, non possono crescere di una quantità superiore a $O(k/\sqrt{N})$ per iterazione, stabilendo un limite inferiore di $\Omega(\sqrt{N})$ interrogazioni.

\section{Confronto con Algoritmi Classici}

\begin{table}[H]
\centering
\begin{tabular}{@{}lcc@{}}
\toprule
\textbf{Approccio} & \textbf{Interrogazioni (caso medio)} & \textbf{Interrogazioni (caso peggiore)} \\
\midrule
Ricerca classica (deterministica)  & $N/2$          & $N$           \\
Ricerca classica (randomizzata)    & $N/2$          & $N$           \\
\textbf{Algoritmo di Grover}       & $O(\sqrt{N})$  & $O(\sqrt{N})$ \\
\bottomrule
\end{tabular}
\caption{Confronto tra ricerca classica e algoritmo di Grover}
\end{table}

Per $N = 10^{12}$: la ricerca classica richiede in media $5 \times 10^{11}$ passi, mentre Grover ne richiede circa $10^6$ -- sei ordini di grandezza in meno.

\end{onehalfspace}
